% ----------------------------------------------------------
% Considerações Finais
% ----------------------------------------------------------
\chapter{Considerações Finais}

O presente trabalho propôs o desenvolvimento de um modelo preditivo de tendência de desempenho acadêmico para estudantes do ensino médio da rede pública estadual do Piauí, fundamentado em fatores sociológicos. Para tanto, foram realizadas duas entregas articuladas: a construção de um conjunto de dados sintéticos com parametrização rastreável a estatísticas oficiais e à literatura acadêmica (no qual os preditores representam sinais parciais do primeiro semestre e o alvo, rotulado pelos critérios de aprovação da LDB, representa a tendência ao final do ano letivo), e a comparação sistemática de três algoritmos de classificação (\textit{Decision Tree}, \textit{Random Forest} e XGBoost) frente a \textit{baselines} de referência, sob o processo KDD, com validação cruzada estratificada ($k=10$), balanceamento por SMOTE confinado a cada \textit{fold} e teste estatístico pareado entre os modelos.

Retomam-se, a seguir, as hipóteses formuladas na Introdução:

\begin{description}
    \item[H1 (confirmada parcialmente).] Os classificadores de \textit{ensemble} superaram a árvore única e os \textit{baselines} com significância estatística (Wilcoxon pareado, $p = 0{,}002$ em ambas as comparações): a \textit{Decision Tree} atingiu 78,22\% de acurácia no \textit{holdout}, a \textit{Random Forest} 82,00\% e o XGBoost 84,11\%, contra 45,00\% do \textit{baseline} majoritário e 80,22\% da Regressão Logística. Entre os dois \textit{ensembles}, contudo, a diferença não é estatisticamente significativa ($p = 0{,}232$): o XGBoost foi adotado como modelo final por sua vantagem numérica consistente, registrando-se a equivalência estatística com a \textit{Random Forest}.
    \item[H2 (confirmada).] Os fatores sociológicos parametrizados no fenômeno gerador foram recuperados pelos modelos como atributos de importância não redundante: após os sinais acadêmicos parciais, as variáveis mais determinantes no XGBoost foram a distorção idade-série, o recebimento do Pé-de-Meia, o trabalho juvenil, a renda \textit{per capita} e o acesso à internet, o que serve como validação metodológica de que a arquitetura captura essa estrutura quando presente nos dados.
    \item[H3 (confirmada).] Os sinais parciais do primeiro semestre foram suficientes para prever a tendência do desfecho anual com desempenho cerca de 39 pontos percentuais acima do \textit{baseline} majoritário, em um cenário com deriva sistemática e idiossincrática de segundo semestre.
    \item[H4 (confirmada).] A regra legal de elegibilidade da assistência financeira produziu, sem qualquer parâmetro dedicado, proporções coerentes de beneficiários entre os cenários: 47,7\% no cenário base, 57,6\% no perfil de alta vulnerabilidade e 6,5\% no perfil de renda alta.
\end{description}

Nos três modelos, nenhum erro ocorreu entre as classes extremas (\textit{Aprovado} e \textit{Reprovado}): todos os equívocos concentraram-se em adjacências ordinais, e a classe crítica \textit{Reprovado} obteve \textit{recall} de 84,3\% no XGBoost (um desempenho relevante para a finalidade preventiva do modelo), pois minimiza os falsos negativos no grupo que mais demanda intervenção pedagógica, em momento do ano letivo no qual a ação preventiva ainda é possível.

O bom comportamento do XGBoost é coerente com suas características arquiteturais: a otimização por aproximação de Taylor de segunda ordem, que incorpora a curvatura da função de perda via Hessiano, e o termo de regularização explícito que penaliza a complexidade das árvores \cite{chen2016}. Essas propriedades permitem ao algoritmo capturar interações não lineares complexas (como a sinergia entre desempenho acadêmico e contexto socioeconômico) que modelos mais simples não representam adequadamente \cite{friedman2001}, ainda que, neste problema, a \textit{Random Forest} tenha alcançado patamar estatisticamente equivalente.

\section{Contribuições do Trabalho}

A principal contribuição desta pesquisa reside na demonstração de que é possível unir, em uma única investigação, a construção de um conjunto de dados sintéticos sociologicamente fundamentado (com rastreabilidade documentada entre cada parâmetro de geração e sua fonte estatística ou acadêmica, incluindo o nível de desagregação e as notas de \textit{proxy}, exportada automaticamente pelo gerador) e a modelagem preditiva comparativa da tendência de desempenho acadêmico. O gerador incorpora mecanismos anti-deterministas ausentes em abordagens ingênuas de simulação: a assistência financeira é derivada de regra de elegibilidade legal (e não de sorteio), as interações entre vulnerabilidades são não lineares, perfis de resiliência acadêmica impedem que os modelos aprendam um determinismo social absoluto, e a rotulagem deriva dos critérios reais de aprovação da LDB, com proporções de classe emergentes.

Adicionalmente, o trabalho contribui para a consolidação de boas práticas metodológicas em Mineração de Dados Educacionais: prevenção de vazamento de dados por meio do SMOTE aplicado exclusivamente no interior de cada partição da validação cruzada \cite{madugula2026}; adoção do F1-Score ponderado como métrica de otimização em cenário desbalanceado \cite{he2009}; inclusão de \textit{baselines} de referência e verificação estatística das diferenças entre modelos \cite{demsar2006}; seleção de variáveis derivadas validada por estudo de ablação; caracterização explícita do teto de acurácia via análise de sensibilidade ao ruído idiossincrático; e salvaguardas de equidade algorítmica, com exclusão de gênero e cor/raça do treinamento. Todo o procedimento é reprodutível e parametrizável, permitindo a simulação de cenários socioeconômicos contrastantes por simples troca de configuração.

\section{Limitações e Implicações}

Não obstante os resultados promissores, algumas limitações devem ser consideradas. A primeira e mais significativa é a natureza sintética da base (em particular, por se tratar de simulação baseada em regras, e não de dados sintéticos derivados de microdados reais): embora cada distribuição marginal seja ancorada em fontes oficiais auditadas, as relações conjuntas entre variáveis refletem as regras especificadas pelo pesquisador. As métricas reportadas medem, portanto, a capacidade dos classificadores de recuperar um fenômeno cuja estrutura (ainda que sociologicamente fundamentada, auditada quanto às fontes e dotada de componente estocástico irredutível) é conhecida por construção, e o desempenho absoluto não deve ser extrapolado diretamente para populações reais. O fenômeno gerador é conhecido por construção, e a transposição para registros reais de estudantes pode revelar padrões, vieses e relações não capturados pela parametrização. A validação com dados reais constitui etapa imprescindível antes de qualquer aplicação prática dos modelos, e as conclusões sobre a importância dos fatores sociológicos devem ser lidas como validação metodológica, não como evidência empírica independente. Como implicação, o valor da evidência produzida está na validação metodológica da arquitetura de classificação, na caracterização do comportamento comparativo dos algoritmos (incluindo a equivalência estatística entre os \textit{ensembles}) e na demonstração de que a estrutura sociológica parametrizada é recuperável pelos modelos a partir dos atributos observáveis.

A segunda limitação refere-se ao recorte do fenômeno: o modelo prevê a tendência do desfecho anual a partir dos sinais do primeiro semestre, sem capturar a dinâmica temporal mais longa da trajetória escolar (transições entre séries, efeitos cumulativos de reprovações sucessivas, sazonalidade da frequência). A terceira concerne à técnica de balanceamento: a literatura recente aponta que a sobreamostragem sintética introduz riscos de privacidade em dados reais (particularmente ataques de inferência de associação), recomendando, para implantações futuras, técnicas de geração com garantias formais de privacidade diferencial \cite{madugula2026, kostopoulos2026}.

Quanto à suficiência dos parâmetros, o conjunto adotado cobre o recorte declarado (sinais parciais intra-ano, contexto socioeconômico individual e desfecho do mesmo ano) e está ancorado nas fontes oficiais e na literatura. Há, contudo, espaço de refinamento: novos atributos de engajamento e de trajetória, substituição do SMOTE pelo SMOTENC nas colunas binárias, repetição da validação cruzada com múltiplas sementes e, sobretudo, a validação com registros reais. Esses refinamentos não invalidam o recorte atual; delimitam a evolução natural do gerador e do procedimento experimental.

Como desdobramentos futuros da pesquisa, destaca-se sobretudo a validação da arquitetura de classificação com registros reais de estudantes (mediante autorização institucional e ética), direção que dialoga diretamente com os trabalhos do grupo de pesquisa voltados à construção de conjuntos de dados institucionais sobre trajetória acadêmica, bem como a integração do modelo a sistemas de gestão escolar, com alertas precoces às equipes pedagógicas, em sinergia com os sistemas \textit{web} de apoio à gestão educacional desenvolvidos no mesmo grupo. Delineiam-se ainda, como extensões metodológicas, a modelagem ordinal e temporal do desfecho, a incorporação de técnicas de IA explicável (SHAP e LIME) para explicações individualizadas das predições, o uso do gerador parametrizável para simulação de contrafactuais de políticas públicas e a adoção de geração sintética com garantias formais de privacidade diferencial nas etapas que envolverem dados reais \cite{kostopoulos2026}.
